\section{Step 40: Request only remaining bounded bytes}
If the running count equals capacity, branch to the supplied private-probe
path. Otherwise subtract the count from capacity. Request the smaller of that
remainder and 4096, using the owned descriptor and destination base plus the
running count. This request is positive because the full-buffer case has already
branched away.

Fill \texttt{io\_request}. Select service 14 and issue one read. A large file
need not cause a single large temporary service allocation; chunking bounds each
request. Both source-buffer cursor arithmetic and the total count use 64-bit
operations. The service's request argument remains a small admitted integer.

\paragraph{Observe} Prepare \texttt{s40-1} and stop before
\texttt{io\_request\_ready}. At the first \texttt{io\_request\_ready}, the normal moving
file has returned 4096 bytes, the requested count was 4096, and the destination
argument named the staging-buffer beginning. The accumulated count is still zero.

\Needspace{7\baselineskip}
\paragraph{Check} Why is this read request never zero?
\begin{itemize}
  \item[A.] A full buffer takes the separate probe path.
  \item[B.] Zero-length reads always report errors.
  \item[C.] The file header guarantees a large file.
\end{itemize}
